\clearpage\SourceLeaf{193}
\begin{center}
{\Large\textsc{Posthumous Papers and Correspondence.}}
\end{center}

\bigskip
\begin{center}[ON THE THEORY OF PARALLELS.]\end{center}
\begin{center}\rule{2.5cm}{0.4pt}\end{center}

\begin{center}
\textsc{Gauss} to \textsc{Taurinus}. Göttingen, 8 November 1824.
\end{center}

\begin{center}Your Honour\end{center}

I read your kind letter of 30 Oct. and the accompanying short essay with some
pleasure, all the more because I am accustomed, in most people who
[make] new attempts at the so-called theory of parallel lines, to find no trace of
a true geometric spirit. I have no other objection (or very little else to object to)
to your attempt than that it is incomplete. Admittedly, your presentation of the proof
that the sum of the three angles of a plane triangle cannot exceed $180^{0}$ still leaves
something to be desired in geometric rigour. But this could be supplied,
and there is no doubt that this impossibility can be proved with the
utmost rigour. The matter stands quite differently with the $2^{\mathrm n}.$ part,
that the sum of the angles cannot be less than $180^{0}$. This is the real
knot, the reef on which everything founders. I suspect that you have not been long
